\chapter{The Field That Acts Where It Isn't}
\label{chap:field-that-acts-where-it-isnt}

Part I built the instrument and read motion off single bodies. Part II sets it among the fields --- and the
field is where the instrument's most distinctive power finally comes into the open. The first lesson the field
teaches is a strange one: it can act where it is not. To the instrument a field's physical content is not its
value at a point but the holonomy of its potential around a loop --- a phase the record picks up by going
around, which survives even where the local field vanishes. And here the chapter divides along the line the
whole book is drawn on. The potential the instrument writes down is ours --- a bookkeeping of the field,
fixable many ways for the same physics --- but the holonomy, the phase a loop actually comes back with, is the
world's, the same under every way we write the potential. The field is the world's reach; the gauge is only
our account of it. And then, in this same chapter, the instrument does something it has not done before. In
Part I it could only \emph{shadow} a loop's residue: Foucault's pendulum let it glimpse a holonomy, but the
continuum formula was a named analogy, not a theorem. Here, in an echo chamber, the instrument \emph{proves}
the residue --- establishes, as a finite theorem about a loop, that a closed path carries a charge of exactly
$\pm 1$. That it carries one, now proved, is the world's; which of the two signs it shows is still the frame's.
This is the moment Part I promised. The chapter opens with the field that acts where it isn't, makes the
instrument's first proof its centerpiece, and closes on energy conservation as the boundary audit of Part I,
written now for a field.

\section{The Aharonov--Bohm Effect: the potential's holonomy}
\label{se:aharonov-bohm}

The sharpest demonstration that a field acts where it is not is the Aharonov--Bohm effect. Split an electron
beam into two coherent branches and send them around either side of a tightly wound solenoid that confines a
magnetic flux $\Phi_B$ to its interior. Outside the solenoid --- everywhere the electrons actually go --- the
magnetic field is zero: $\mathbf{B} = 0$ along both paths, so by the classical reckoning there is no force on
either branch and nothing should happen. Yet when the branches recombine, they interfere with a measurable
phase difference, and that difference depends on the flux the electrons never touched.

It is worth pausing on how strange this is. Nothing the electron passes through exerts a force; the field is
sealed inside the solenoid, and the electron's whole journey is through empty space. And yet the two branches,
recombining, disagree by an amount set precisely by the flux they went around but never entered. The
instrument is not detecting a field along the path --- there is none to detect --- it is reading a property of
the loop the two branches enclose between them. The field's reach is not its local value but the residue it
leaves on any loop that encircles it.

The instrument reads the phase as a holonomy --- a residue the loop carries. Each branch transports a phase
along the connection $\mathbf{A}$, the potential we wrote to describe the field, accumulating it step by step,
and the two branches, enclosing the solenoid between them, differ by the loop integral of that potential,
\begin{equation}
  \Delta\varphi = \frac{e}{\hbar}\oint \mathbf{A}\cdot d\boldsymbol{\ell} = \frac{e}{\hbar}\,\Phi_B .
\end{equation}
This integral is not a number ground out of a field along the path --- the field is zero there. It is the
\emph{name of what the loop encircles}: the holonomy of the potential, which equals the enclosed flux exactly,
whether or not anything along the way pushed the electron. As the confined flux is varied, the phase marches
through full turns, period by period in units of the flux quantum $h/e$, a staircase the instrument reads off
the interference and nothing else.

Two facts about the potential make the holonomy the right object, and both reach forward into the quantum
part. The first is gauge invariance. The potential is not unique: one can add to it the gradient of any
function,
\begin{equation}
  A_\mu \to A_\mu + \partial_\mu \chi,
\end{equation}
without changing a single physical prediction, because the added gradient integrates to zero around any closed
loop. The choice of $\chi$ is ours --- a whole family of potentials writes the same field, and which one we use
is bookkeeping --- but the holonomy $\oint \mathbf{A}\cdot d\boldsymbol{\ell}$ is blind to that choice: it sees
only the enclosed flux, the gauge-invariant content the loop carries, the part that is the world's and not ours
to set. It is the line the whole book is drawn on, sharpened to a case: the gauge is the description, the
holonomy is the described. The second fact is what the flux is, locally:
the curl of the potential, the field tensor
\begin{equation}
  F_{\mu\nu} = \partial_\mu A_\nu - \partial_\nu A_\mu,
\end{equation}
whose magnetic part is the $\mathbf{B}$ sealed inside the solenoid. Along the electron's path $F_{\mu\nu} = 0$
--- the field is zero --- and yet the loop integral of $A$ does not vanish, because the flux $F_{\mu\nu}$
encloses is nonzero even where $F_{\mu\nu}$ itself is not. The potential carries, around the loop, what the
field hides inside it.

And here the forward-link must be stated plainly, because it is the second great through-line the front half
lays into the quantum part. The field tensor $F_{\mu\nu}$ is the abelian precursor of the object the proved
heart of the book is built on: the curvature of a covariant derivative, written there as a commutator,
\begin{equation}
  [D_\mu, D_\nu] = iq F_{\mu\nu} .
\end{equation}
The relation is exactly the one the force chapter set up between the Poisson bracket and the commutator. As
the bracket $\{q, p\} = 1$ is the classical shadow of $[x, p] = i\hbar$, so the field tensor here --- the plain
curl of a potential that commutes with itself --- is the abelian, gentle version of the non-abelian field
strength the Dirac chapter carries, where the potentials no longer commute and the curvature gains a term
built from their bracket. The holonomy the electron picks up here is the same loop residue that, sharpened, the
later chapters prove to be $\pm 1$.

The honest floor is a finite count obligation: the phase is a count the instrument accumulates around a loop,
required to equal the enclosed flux, and that requirement is exact --- the holonomy is not approached by
refinement but fixed by what the loop encircles. The counting is ours; what the count is required to equal is
the world's. The reading is that the potential, not the local field, is what the instrument carries. A field acts where it is not through the holonomy of its connection: the record, going
around a loop, picks up what the loop encircles even when nothing along the path was there to push it. The
loop residue of the previous part has become a physical phase --- and the next effect makes it a proof.

\section{The Echo Chamber Maze: the first proved holonomy}
\label{se:echo-chamber-maze}

Here is the landmark, and it deserves to be marked as one: the first time in the book the instrument does not
shadow a loop's residue but \emph{proves} it. Picture navigating a maze by clapping and listening, so the
echoes trace the causal paths a signal takes. Straight corridors --- a flat metric --- return crisp echoes,
perfect parallel transport with no phase; curved passages distort the return, leaving a residue. Walk a closed
loop through the maze and compare the echoed rhythm against what was sent, and the mismatch measures the
curvature the loop encloses. An open path carries no residue. A closed loop carries $\pm 1$ --- that it carries a sign at all is the
world's, forced by the curvature the loop encloses; which of the two signs it shows is the frame's, set by the
orientation the walk gives the loop.

The residue is the discrete echo of the last effect's continuous phase. Where the Aharonov--Bohm loop carried
the holonomy $\oint \mathbf{A}$ as a phase that marched smoothly with the enclosed flux, the maze carries a
discrete curvature, and the residue the echo accumulates is the quantized version of that same loop integral:
the enclosed curvature is no longer a continuum the phase sweeps through but a charge the loop is forced to
carry, exactly $\pm 1$. The continuous holonomy and the discrete residue are one object at two
resolutions --- and it is the discrete one, here, that the instrument can prove.

What makes this different from everything before is the \emph{status} of that $\pm 1$. When Foucault's
pendulum traced its loop in Part I, the instrument could only shadow the result: the discrete failure-to-close
approached a continuum holonomy, and the bridge between them was an analogy, honestly named and no more. Here
the instrument crosses the line it has stood behind for four chapters. It does not approach the residue by
refinement; it establishes it, as a finite theorem about the loop --- that an open path is trivial and a
closed loop is charged, exactly $\pm 1$, with the converses to match: off the loop, nothing; around it, the
sign. That the charge is there, and is a unit, the instrument now proves; which way the unit points stays the
frame's to set. The residue Foucault let the instrument glimpse is now a thing it can prove.

To see why this is a proof and not a measurement, watch the instrument check it both ways. Send a clap down a
straight corridor and back, and the echo returns in step --- residue zero, exactly, not approximately. Send it
around a closed loop that encircles a curved cell of the maze, and the echo returns shifted by the full
residue --- $\pm 1$, exactly, again with no margin. The instrument does not refine its way toward these
values; it establishes them as the only admissible answers, the open path forced to nothing and the closed
loop forced to the sign, each converse closing the case. A shadow would say ``the residue is approximately
$\pm 1$, and closer as the maze is read more finely''; a proof says ``the residue is $\pm 1$, and an open path
is exactly trivial,'' and means it. That gap --- between approached and established --- is the gap this
chapter crosses, and every later proof in the book widens the same crossing.

This is the strongest grounding the book has, and it will recur: every later proof --- a diffracting crystal,
a rotating ring, a Dirac particle, a positron's annihilation --- is this same move, a loop residue established
as a finite theorem rather than approached as a shadow. The honest floor is a finite count obligation, and a
proved one: a finite theorem about a loop's residue, not an analogy to a continuum. The proof is described
here, not performed on the page; what matters is its character. The reading is that curvature is the
informational stress of maintaining closure in a finite domain, read off a loop and \emph{proved} to be
$\pm 1$ --- the unit the world's to charge, the sign the frame's to name. The field's geometry, here, is not
approached by refinement; it is demonstrated --- and the instrument has shown, for the first time, that it can
prove and not merely shadow.

\section{The Conservation of Energy: the boundary audit for a field}
\label{se:conservation-of-energy}

The chapter closes where Part I's momentum chapter left a promise. There the instrument audited a boundary: it
drew a closed line around a region --- ours to place where we like --- tallied what crossed in and out, and
read conservation off the books coming out even. The line is ours; that the books come out even is not. Energy
conservation, read off a field, is exactly that audit, raised from a single body's momentum to the full
stress-energy of a field. Take a field $\phi$ with Lagrangian
\begin{equation}
  \mathcal{L} = \tfrac{1}{2}\,\partial_\mu\phi\,\partial^\mu\phi - \tfrac{1}{2}\,m^2\phi^2,
  \qquad \eta_{\mu\nu} = \mathrm{diag}(-,+,+,+),
\end{equation}
and form its symmetric stress--energy tensor,
\begin{equation}
  T^{\mu\nu} = \partial^\mu\phi\,\partial^\nu\phi - \eta^{\mu\nu}\mathcal{L} .
\end{equation}
This tensor is the conserved current Noether's theorem attaches to the field's symmetries --- its energy
component the count that time-translation invariance conserves, exactly as the least-principle chapter
promised, now carried from a single body's motion to a field that fills space. It is the field's complete
ledger of what can cross a boundary: its time--time component is the
energy density the region holds, and its time--space components are the energy flux that crosses the edge,
\begin{equation}
  \mathcal{E} = T^{00} = \tfrac{1}{2}\bigl(\dot\phi^2 + |\nabla\phi|^2 + m^2\phi^2\bigr),
  \qquad
  S^i = T^{0i} = \dot\phi\,\partial^i\phi .
\end{equation}
The conservation law is the statement that this ledger does not leak,
\begin{equation}
  \partial_\mu T^{\mu\nu} = 0,
\end{equation}
which for the energy component is a continuity equation, $\partial_t \mathcal{E} + \nabla\cdot\mathbf{S} = 0$,
and in integral form is the boundary audit itself,
\begin{equation}
  \frac{d}{dt}\int_R \mathcal{E}\, dV = -\oint_{\partial R} \mathbf{S}\cdot\mathbf{n}\, dA .
\end{equation}
What a region loses in energy, it loses through its boundary, accounted for by the flux --- the same
closed-line tally the instrument ran for momentum, now keeping the field's books. Where we draw the boundary is
ours; that the flux across it accounts for every unit the interior lost is the world's, a debt it collects at
whatever edge we draw.

The honest floor is a finite count obligation: the continuity is a count balanced across a boundary, the
field-theoretic form of the momentum balance that closed the mechanics part. The reading is that energy is not
a substance a region stores but a count that balances across its boundary --- the audit ours to run, the
balance the world's to enforce --- conserved by the same continuity that conserved momentum, Part I's
bookkeeping written now for a field, and audited at the edge exactly as before.

\bigskip

Part II has opened on the instrument's two new powers. A field acts where it is not, through the holonomy its
potential carries around a loop --- the potential ours to write, the holonomy the world's to enforce; and the
instrument, for the first time, has proved such a residue rather than shadowed it --- the line Foucault let it
approach, now crossed, with the sign it carries still the frame's to set. The boundary audit of Part I has
carried over intact, keeping the field's energy books at the edge --- the audit ours, the balance the world's.
The next chapter stays with the field but turns to what it does when it overlaps itself --- diffraction,
interference, and the wall of crossed polarizers --- where the instrument meets its second no-go and the proved
residue of this chapter is put to work reading patterns of light.

\bigskip
\begin{center}
\rule{0.4\textwidth}{0.4pt}
\end{center}
\bigskip

\noindent The most telling evidence in a case is often for something no one on the scene ever saw. A cause can
act where it is not --- sealed away, untouched by any path the record ran through --- and still leave its mark,
an effect entered plainly in the account though the thing behind it was never in reach. The instrument admits
such a cause the only way it honestly can: not by witnessing it, but by a trace the record could not otherwise
carry. What no one on the path could see stands in the account all the same, on the strength of what only that
cause explains. The account is ours to keep; the mark it is forced to carry, by something it never touched, is
the world's.
